\chapter{Introduction}
\label{ch:intro}

A supernumerary robotic limb is an extra arm worn on the body, working
alongside the wearer's own limbs rather than replacing or rehabilitating
them~\cite{yang2021srlreview,prattichizzo2021augmentation}. What separates
such a device from a fixed manipulator is consequential for control: the base
of the robot is a person. It accelerates, sways and steps, and it cannot be
told to hold still.

This thesis reports work on the MUlti-limb Virtual Environment, known within
Imperial College London as the Dr.\ Octopus platform~\cite{muve}, which mounts
two Kinova Gen3 seven-degree-of-freedom arms~\cite{kinova2022gen3} on a
backpack frame and commands them from a \emph{second} person: an operator who
cannot feel the arms, controlling them for a wearer who cannot command them.
Fusion established that this operator/wearer architecture can be built, but
measured nothing about it~\cite{saraiji2018fusion}; shared autonomy is known
to raise operator performance where the robot is not attached to
anybody~\cite{dragan2013policyblending,javdani2015hindsight}; and increasing
apparent autonomy has been shown to lower a wearer's trust in a setting where
autonomy was simulated by a concealed human with no task of their
own~\cite{zhou2026proxemics}. No study measures the operator and the wearer
together, under the same autonomy, in the same trial, so the exchange rate
between what shared autonomy buys the operator and what it costs the wearer
is unknown. \Cref{ch:background} reviews this ground in full; the gap it
leaves is what this project is built to close.

Two premises, argued and tested throughout, shaped the engineering. First,
commanding fourteen joints through a body-worn master with no force feedback
is a usability problem before it is a control problem, so position is taken
from whichever sensor observes each component best rather than from the
slave's full kinematic chain. Second, a body-worn master \emph{will} lose
sensing channels over its life, so the system treats channel loss as an
operating condition: every channel carries a health verdict, and when too
many are lost the system reports position unavailable rather than fabricating
one, because on a manipulator beside a person's head a fabricated position is
worse than none.

\section{Objectives}

\begin{enumerate}
  \item Design and build an instrumented mannequin master reproducing the
    Gen3 morphology at reduced scale, characterised against recorded data.
  \item Derive a master-to-slave command mapping that degrades predictably as
    sensing channels are lost, and quantify what each loss costs.
  \item Implement a safety architecture appropriate to a manipulator mounted
    beside a person, validated by fault injection rather than inspection.
  \item Characterise the reachable workspace of the mounted configuration and
    determine which manipulation tasks it can support.
  \item Establish a verification methodology producing traceable evidence for
    every claim, and apply it throughout.
  \item Run a pilot comparison of direct control against shared-autonomy
    assistance with real operators, and identify what stands between it and
    the full two-person study the architecture was designed for.
\end{enumerate}

The planning report framed this project around a haptic master with physical
springs and virtual slave-side impedance. That objective was not delivered:
\cref{ch:discussionmain} explains why, and \cref{app:development} records the
route in full.

\section{Contributions}

Each contribution below is stated with its evidence class, because the
distinction between implemented, verified in simulation, measured on
hardware, and measured on a pilot session is load-bearing throughout this
report. A graceful-degradation architecture for a body-worn master, measured
against \num{20430} recorded frames; a reduced-degree-of-freedom command
mapping that takes direction from gravity and magnitude from the
best-conditioned joints; a safety architecture in which every blocking
mechanism is observable, with fourteen of fourteen injected faults handled; a
workspace characterisation showing the two arms share no reachable point
under the as-built mount, together with a quantified and buildable mount
correction; a jerk-limited, phase-synchronised motion generator that removes
\SI{51.3}{\milli\metre} of unrequested Cartesian excursion; a
calibrate-then-plan perception stage that recovers six objects to within
\SI{1}{\milli\metre}; a stage-by-stage instrumentation of the perception
pipeline exercised on four physical cameras; and a first pilot of the
direct-versus-shared-autonomy comparison with four operators across three
tasks, run under ethical approval once VR controllers supplied a working
alternative to the master's incoherent channels. \textbf{No contribution
below the pilot was validated on physical hardware carrying a load.}

The single most transferable outcome does not appear in that list. Twenty-six
times across this project, a result that looked like a system failure was a
defect in the measuring tool rather than in the system, and validating an
instrument against a known answer before trusting a surprising reading it
produces is the practice applied to every number in this report
(\cref{sec:validity}, \cref{ch:instrument}).

\section{Report structure}

\Cref{ch:methodsmain} summarises what was designed, built and run, including
the pilot protocol and its ethical approval; the full engineering detail --
the master's electronics and firmware, the system architecture, the
perception and motion pipelines -- is in \cref{ch:master,ch:method,ch:perception,ch:motion}.
\Cref{ch:resultsmain} reports the headline measurements and the pilot's
results; \cref{ch:results,ch:vision,ch:workspace} carry the full simulation
campaign, the real-camera vision session and the workspace characterisation.
\Cref{ch:discussionmain} interprets the results and states what they are
worth; \cref{ch:discussion} is the extended version. \Cref{ch:conclusion}
closes against the objectives above. \Cref{app:study,app:analysis} are the
protocol for the full two-person study and the validated analysis that would
evaluate it, and \cref{app:pilot} is the pilot session in full.
